\textbf{Proof.}
Let \(C=\binom D2\).  Each unordered vertex pair has at most three basal states in \(E\): absent, oriented in one direction, or oriented in the other direction.  Thus at most \(3^C\) directed graphs are visited before cyclic choices are rejected.  There are \(2^D\) subsets \(A\) before checking \(\operatorname{an}_E(A)=A\).  Once \((E,A)\) is fixed, \(\texttt{def:normalized-dai-certificate}\) uniquely computes
\[
Q_{\mathbf M}(E,A)=S_0(\mathbf M)\setminus B(E,A),
\]
and, for each \(k\in\mathcal K_+\), uniquely computes
\[
Z_k(\mathbf M),\quad R_k(\mathbf M,A),\quad O_k(\mathbf M,A),
\quad T^*_{k,\mathbf M}(E,A),\quad F^*_{k,\mathbf M}(E,A).
\]
The residual-skeleton and descendant-closure tests either reject the pair or admit that single derived certificate; they introduce no choices.  Therefore the enumeration visits at most
\[
3^C2^D=3^{\binom D2}2^D
\]
candidates and terminates.  If \(K=0\), all displayed nonobservational objects are absent and this count is unchanged.

Let
\[
\mathfrak A(\mathbf M)
=\left\{(E,A):
\operatorname{Check}_{D}(\mathcal R_{\mathbf M}(E,A),\mathbf M)=1
\right\}.
\]
Assumption \(\texttt{ass:compatible-input}\) and \(\texttt{thm:finite-common-realizability}\) imply \(\mathfrak A(\mathbf M)\ne\varnothing\).  Fix \(e=(k,u;uv)\in\mathcal E(\mathbf M)\).  The two directions and exhaustiveness clause of that theorem yield
\[
\left\{
\operatorname{mark}_e\bigl(\mathcal N^{(k)}(\mathcal R_{\mathbf M}(E,A))\bigr):
(E,A)\in\mathfrak A(\mathbf M)
\right\}
=
\left\{
\operatorname{mark}_e\bigl(\mathcal M_\eta^{(k)}\bigr):
\eta\in\mathfrak C_D(\mathbf M)
\right\}.
\]
For the left-to-right inclusion, every passing pair reconstructs a compatible member.  For the reverse inclusion, target residualization of every compatible member supplies a passing pair with the same entire fully oriented regime-MAG family.

Every fully oriented MAG endpoint is a tail or an arrowhead.  Hence the common feasible set is one of
\[
\{\mathrm{tail}\},
\qquad
\{\mathrm{arrow}\},
\qquad
\{\mathrm{tail},\mathrm{arrow}\}.
\]
The output rule in \(\texttt{def:complete-cdis}\) returns respectively a tail, an arrowhead, or a circle, exactly as \(\texttt{def:canonical-si-pag}\) does.  The accepted certificates also have the supplied skeleton in each regime.  Consequently, endpoint by endpoint,
\[
\operatorname{CDIS}^{\mathrm{cmp}}(\mathbf M)=\mathcal P_{\mathrm{SI}}(\mathbf M).
\]

If \(e\) is circular, its feasible set contains both marks.  The finite lexicographic enumeration therefore retains one accepted tail certificate and one accepted arrowhead certificate.  By \(\texttt{thm:finite-common-realizability}\), their canonical reconstructions are actual members of the common-basal class \(\mathfrak C_D(\mathbf M)\).  At every fixed output endpoint the feasible set is a singleton, so both witnesses agree with every fixed mark; at \(e\) they realize opposite marks.  These are coupled whole-family witnesses, not independent regime-wise MAG completions.

For the runtime, materializing \(S_0(\mathbf M)\) and all indicator-neighborhoods costs \(O(D^2+KD)\).  For a candidate \((E,A)\), acyclicity, transitive closure, ancestor closure, all descendant closures, and all minimal-generator sets are computable in
\[
O((K+1)(D+1)^3)
\]
time.  Constructing \(B(E,A)\) checks at most \(D\) possible common children for each of \(O(D^2)\) pairs and is absorbed by the same cubic envelope.  Forming the derived residuals, constructing all regime matrices, deciding the path conditions, and comparing all unshielded triples are likewise covered by that envelope.  Updating endpointwise feasible sets and retaining the first opposing certificates costs \(O(N_{\mathrm{end}})\).  Thus the conservative per-candidate cost is
\[
O\!\left(D^2+(K+1)(D+1)^3+N_{\mathrm{end}}\right).
\]
Multiplication by the candidate count gives
\[
O\!\left(3^{\binom D2}2^D
\left[D^2+(K+1)(D+1)^3+N_{\mathrm{end}}\right]\right).
\]
The logarithm of the search factor is \(O(D^2)\); \(K\) appears only in the displayed polynomial factor.  This is a non-tight singly exponential upper bound in \(D^2\), not a polynomial algorithm.  It proves no completeness result for a fixed local orientation calculus and excludes the conditional-independence and invariance queries required to obtain the supplied oracle MAGs.  No support, overlap, rank, inverse, division, or limiting condition is used. \(\square\)